\section{Conclusion}
\label{sec:conclusion}

This paper compares 16 frequency estimators over 34 scenarios in a local
relay-class measurement setting that combines standard dynamic tests with
IBR-specific composite stress. The results show disturbance dependence:
estimator rankings move with both the disturbance family and the metric used to
read them. Standard compliance-style tests remain necessary reference cases,
but they are not sufficient by themselves for protection-oriented evaluation in
IBR-dominated grids.

We highlight three findings. First, the latency--robustness trade-off is
structural: nine estimators win at least one RMSE scenario, and no estimator
wins more than 10 of the 34. Windowed spectral methods reduce steady harmonic
error but pay observation delay, while low-cost loop-based methods remain
attractive under tight computational budgets. Second, Kalman-family methods,
especially EKF and RA-EKF in this dataset, remain competitive across the
widest range of families. RA-EKF is evaluated here as one Kalman-family
variant rather than as a new estimator family. Third, the hardest cases are
the IBR multi-event sequence, the $60^\circ$ phase-jump cases, and the noisy
ringdown variants; these are also the cases in which RMSE, peak error, and
trip-risk select different estimators.

The study is limited to the scenario set and relative computational-cost
measurements reported here. CPU time is used as a software-level cost indicator,
not as hardware latency, and the study remains simulation-only. Hardware-in-the-loop
validation, wider coverage of converter-control interactions, and sensitivity
analysis of the tuning protocol are the natural next steps.

From a deployment standpoint, the benchmark identifies a small practical core.
SOGI-PLL is the strongest low-cost choice for ramp-dominated operation. EKF is
the most consistent model-based baseline for magnitude steps, out-of-band
interference, and balanced composite-event behavior. RA-EKF earns its place not
as a universal winner but as a higher-cost Pareto point with the best mean RMSE
and the strongest modulation-family results. IpDFT and TFT remain relevant when
limiting trip-risk matters more than minimizing average error, provided their
observation delay is acceptable.

In practice, the benchmark does not support a single default estimator for all
low-inertia IBR conditions. It supports a screening workflow: narrow
candidates by disturbance family and compute budget, then resolve the final
choice with the metric aligned to the protection objective. In that workflow,
ramps mostly test tracking bandwidth, jumps test recovery after
discontinuities, and harmonic-rich cases test robustness to leakage and
interharmonics.
